% ***************************************************
% Chapter 6: Conclusion and Future Roadmap
% ***************************************************

\chapter{Conclusion and Future Roadmap}
\label{chap:conclusion}

\section{Summary of Research Progress and Achievements}
\label{sec:conclusion_summary}
This progress report has detailed the conception, technical implementation, and empirical evaluation of primitive-space scene change detection applied to underwater 3D Gaussian Splatting models. In the initial phase of the research, all primary computational milestones have been met:
\begin{enumerate}
    \item \textbf{Automated Cluster Infrastructure:} A modular, production-grade 6-stage automation suite was designed and deployed across 11 isolated Conda environments on the QUT Aqua High-Performance Computing cluster, eliminating manual intervention in large-scale multi-view photogrammetry.
    \item \textbf{Master 15 / 15 Benchmark Milestone:} All 15 planned benchmark targets across 8 models and 3 datasets achieved 100\% full-lifecycle completion, demonstrating rock-solid execution stability on NVIDIA A100-SXM4-40GB hardware.
    \item \textbf{Comprehensive Model Benchmarking:} Quantitative evaluation across held-out novel views revealed critical architectural trade-offs:
    \begin{itemize}
        \item 3D-UIR demonstrated the highest structural preservation ($SSIM = 0.8739$ on Kwaj, $0.7085$ on Tokai), proving optimal for retaining fine geometric coral features without blurring.
        \item SeaSplat successfully decoupled water column backscatter and wavelength attenuation ($PSNR = 24.60\text{ dB}$, $SSIM = 0.8341$), recovering true benthic albedo.
        \item Gaussian Splashing achieved the highest raw reconstruction fidelity ($PSNR = 27.93\text{ dB}$ on Kwaj), while WaterSplatting achieved the lowest perceptual error ($LPIPS = 0.1714$) with sub-40-minute training times.
    \end{itemize}
    \item \textbf{Per-Model Surface Regularization and 3D Mesh Review (SuGaR):} Motivated by the recommendation at the conclusion of the primitive-space change detection paper~\cite{galappaththige2026pixelsprimitivesscenechange}, enforced signed distance field (\gls{sdf}) priors to prune 85--92\% of floating water-column primitives and engineered per-model output directory isolation so that every evaluated underwater method exports its own distinct, watertight textured 3D surface mesh (OBJ format) for comparative review alongside the primary 3D Gaussian splats.
    \item \textbf{Primitive-Space Change Detection (O-SCD):} Validated the primary primitive-space differencing engine on the Custom OSCD dataset, achieving a PSNR of 27.08~dB and SSIM of 0.8315 while accurately separating added and removed objects directly in 3D Gaussian space without viewpoint dependencies.
\end{enumerate}

\section{Technical Challenges and Lessons Learned}
\label{sec:lessons_learned}
Experimental execution yielded valuable insights that inform subsequent project phases:
\begin{itemize}
    \item \textbf{Gaussian Drift and SuGaR Surface Anchoring:} Unconstrained Gaussian splats optimize purely for photometric re-projection; independent runs on identical scenes exhibit spatial centroid drift. Integrating SuGaR surface regularization, as recommended at the conclusion of the 3D change detection literature~\cite{galappaththige2026pixelsprimitivesscenechange}, and reviewing the resulting 3D OBJ meshes across underwater methods proves indispensable for anchoring primitives to physical geometry and diagnosing whether higher-fidelity 3D reconstructions can resolve fine-scale $\text{mm} \to \text{cm}$ coral changes.
    \item \textbf{Per-Model Mesh Cache Isolation:} Vanilla SuGaR defaults to scene-level output paths and skips coarse extraction if a file already exists, causing multiple models on the same scene to share identical meshes unless isolated via model-specific directory tagging.
    \item \textbf{Severe Turbidity and Long-Range Collapse:} In highly turbid environments (such as Tokai Maru), optical attenuation models require accurate metric depth priors to avoid collapsing background Gaussians into foreground haze.
    \item \textbf{Batch Script Robustness:} Diverse submodules with conflicting GCC/CUDA dependencies require rigid environment sandboxing and process isolation.
\end{itemize}

\section{Future Roadmap: The AIMS EcoRRAP Deployment}
\label{sec:future_roadmap}
With baseline benchmarks and cluster toolchains fully validated, the project is strategically positioned to enter Phase~2 and Phase~3 upon receipt of the 1.35~TB AIMS EcoRRAP dataset:

\subsection{Phase 2: EcoRRAP Ingestion and Co-Registration (Months 1–2)}
\begin{itemize}
    \item \textbf{Data Ingestion \& Staging:} Finalize the inter-institutional Data Transfer Agreement (DTA) with AIMS, staging the 1.35~TB temporal dataset onto the QUT Aqua Lustre scratch filesystem via high-throughput Globus transfer.
    \item \textbf{Multi-Temporal SfM Alignment:} Ingest multi-temporal image pairs ($T_0$ baseline vs. $T_1$ post-disturbance) through the Stage~2 COLMAP pipeline, establishing unified extrinsic camera poses across seasonal surveys.
    \item \textbf{Depth Prior Extraction:} Run batch monocular depth estimation across the coral scenes to provide dense spatial constraints for optical inversion.
\end{itemize}

\subsection{Phase 3: Model Transplantation and Extended Cluster Optimization (Months 3–4)}
Figure~\ref{fig:ecorrap_architecture} illustrates the target Phase~3 architectural integration bridging underwater optical restoration, \gls{sugar} surface regularization, and primitive-space scene change detection.
\begin{itemize}
    \item \textbf{Architectural Transplantation:} Based on the benchmark findings and per-model 3D surface mesh evaluations, directly graft the winning underwater optical formulations (3D-UIR and SeaSplat) into the \gls{oscd} pipeline, replacing the default naive terrestrial Gaussian backbone.
    \item \textbf{Long-Duration HPC Training:} Train the unified model for extended iteration counts (50,000 to 100,000 iterations) on multi-GPU A100/H100 nodes, generating ultra-dense, turbidity-free representations of baseline reef colonies.
    \item \textbf{Millimetre-to-Centimetre ($\text{mm} \to \text{cm}$) Coral Change Quantitation:} Execute primitive-space change detection on the 3D Gaussian splats against the $T_1$ inspection surveys, supported by per-model SuGaR 3D mesh extraction to resolve fine-scale $\text{mm} \to \text{cm}$ coral changes and calculate exact physical metrics:
    \begin{itemize}
        \item Millimetre-to-centimetre ($\text{mm} \to \text{cm}$) skeletal branch extension and volumetric accretion ($\text{cm}^3$) indicating coral growth and calcification.
        \item Volumetric depletion ($\text{cm}^3$) reflecting cyclone breakage, storm damage, or coral mortality.
        \item Surface area changes ($\text{m}^2$) quantifying structural rugosity shifts.
    \end{itemize}
\end{itemize}

\begin{figure}[htbp]
    \centering
    \begin{tikzpicture}[
        node distance=0.55cm and 0.8cm,
        block/.style={rectangle, draw=blue!75!black, fill=blue!6, thick, text width=4.4cm, minimum height=1.15cm, align=center, font=\scriptsize\sffamily, inner sep=3pt},
        wideblock/.style={rectangle, draw=blue!80!black, fill=blue!10, thick, text width=6.8cm, minimum height=1.2cm, align=center, font=\scriptsize\sffamily, inner sep=3pt},
        arrow/.style={->, >=stealth, thick, blue!75!black}
    ]
        \node[block] (t0) {\textbf{Baseline Survey ($T_0$)}\\[1pt]Raw EcoRRAP Coral Imagery};
        \node[block, right=1.4cm of t0] (t1) {\textbf{Inspection Survey ($T_1$)}\\[1pt]Post-Disturbance Reef Imagery};
        \node[block, below=of t0] (uw) {\textbf{Underwater-Optimised 3DGS}\\[1pt]3D-UIR / SeaSplat Optical Inversion};
        \node[block, below=of uw] (sugar) {\textbf{SuGaR Surface Regularization}\\[1pt]Enforce SDF Prior \& Prune Floaters};
        \node[wideblock, below=3.4cm of $(t0)!0.5!(t1)$] (oscd) {\textbf{Primitive-Space O-SCD Engine}\\[1pt]Turbidity-Invariant 3D Gaussian Differencing ($T_1$ vs.\ $T_0$ Prior)};
        \node[wideblock, below=of oscd] (vol) {\textbf{Differential SuGaR Mesh \& Ecological Metrics}\\[1pt]Volumetric Accretion/Depletion ($\text{cm}^3$) \& Rugosity ($\text{m}^2$)};

        \draw[arrow] (t0) -- (uw);
        \draw[arrow] (uw) -- (sugar);
        \draw[arrow] (sugar.south) |- ([yshift=0.25cm]oscd.north) -- (oscd.north);
        \draw[arrow] (t1.south) |- ([yshift=0.25cm]oscd.north) -- (oscd.north);
        \draw[arrow] (oscd) -- (vol);
    \end{tikzpicture}
    \caption{Target Phase~3 unified temporal underwater change detection architecture, combining underwater-optimised 3DGS optical restoration and SuGaR surface regularization with primitive-space O-SCD for the 1.35~TB AIMS EcoRRAP dataset.}
    \label{fig:ecorrap_architecture}
\end{figure}

\subsection{Phase 4: Synthesis, Validation, and Handover (Month 5)}
\begin{itemize}
    \item \textbf{Ecological Validation:} Review volumetric change maps with domain experts (Dr. Scarlett Raine and AIMS researchers) to validate biological plausibility against in-situ diver transects.
    \item \textbf{Dissemination:} Compile the final Master's thesis, draft a conference manuscript for submission to a leading vision venue (e.g., IEEE ICRA or IROS), and package the codebase for long-term deployment at the QUT Centre for Robotics.
\end{itemize}
